\section{Outlook: Macroeconomics}

\begin{frame}{What we have not dealt with yet}
    \mytext{
    \begin{itemize}
        \item money, credit, interest rates
        \item the economy as a whole: macroeconomics
        \item output (gross domestic product)
        \item savings (in the economy as a whole)
        \item investment (in the economy as a whole)
        \item consumption (in the economy as a whole)
        \item international economic relations
    \end{itemize}
    }
\end{frame}

\begin{frame}{Circular Model of the Economy}
    \showmytikz[1,2,3,4,5,6]{../shared/tikz/tikz-flowmodel-own}
    \mycomment{
        In the context of macroeconomics, the economy is circular: households save (on capital markets), their savings become investments (of business). Households buy and consume (goods and services), which are produced by businesses. 
        
    }
\end{frame}

\begin{frame}{Circular Model of the Economy}
    \showmytikz[1,2,3,4,5,6]{../shared/tikz/tikz-flowmodel}
    \mycomment{
        Households pay taxes (to the government), which the government spends on markets (to businesses). Households import goods and services from the rest of the world, while business export to the rest of the world.

        Note that this is simplified. The government may also redistribute, and has other functions as well. Households decide how much time (and skill) to offer on labour markets. There is also an international capital market. Businesses may be organised in one way or the other, and households may be large or small.
    }
\end{frame}

\begin{frame}{Value and Money}
    \mytext{
        There is a clear distinction between
        \begin{itemize}
            \item use value
            \item exchange value
        \end{itemize}
        of any good.
    
        The exchange value is expressed by the price of the good on the market.
    
        The use value is subject to the preference of the owner.
    }
    \mycomment{
        If goods are tradable, every good will have at least the exchange value for everyone; while some people may assign a higher use value to the good. The latter means that they will not part from the good at exchange value; if my personal evaluation of my car is larger than its retail value, I will keep it - for what we may call sentimental reasons.
    }
\end{frame}

\begin{frame}{Value and Money}
    \mytext{
    What is the value of money?
    
    The four functions of money are: 
    \begin{enumerate}
        \item medium of exchange
        \item unit of account
        \item store of value
        \item standard of deferred payment
    \end{enumerate}
    % \begin{itemize}
    %     \item role of banks and central bank?
    %     \item different currencies, etc.?
    %     \item not really a topic in micro-economics
    %     \item real income vs. nominal income
    % \end{itemize}    
    }
    \mycomment{
    The state is the institution that issues currency. This is  the descretion of the (independent) central banks. The effects of monetary policy is studied in macroeconomics.
    }
\end{frame}

\begin{frame}{Value and Money}
    \mycomment{
    The exchange value of a good is a result from supply and demand; costs play indirectly a role via supply.

    Imagine two countries, each of which produces only one good: one coconuts, one fish. They swap coconut for fish depending on their preferences; there will be some exchange rate fish per coconut.

    There is a market fish-coconut only.
    }
    \showmytikz[1]{../shared/tikz/tikz-quantitytheoryofmoney}
    
\end{frame}

\begin{frame}{Value and Money}
    \mycomment{
    
    Now imagine they start using some sort of extinct shells as means of exchange or for deferred payment.
    
    There is now a market fish-for-shells, and a market coconut-for-shells.
    
    What will the price of fish and coconut depend on?

    It will depend on the quantity of shells, and on the speed of circulation of shells. 
    
    This is the so-called "quantity theory of money": the value of money results plainly of its quantity.
    }

    \showmytikz[2]{../shared/tikz/tikz-quantitytheoryofmoney}
    
\end{frame}


\begin{frame}{Exchange Rates: Purchasing Power Parity}
    \mytext{
        Exchange rate assures the same basket of goods can be baught here or there

    (`Big Mac Index')

    \begin{itemize}
        \item Often used to `correct' exchange rates (space)
        \item often used to correct value comparisons over time
    \end{itemize}
    }
\end{frame}

\begin{frame}{Exchange Rates: Purchasing Power Parity}
    \mytext{
      Price of haircut in Switzerland 80 SFR $\sim 90 $ \$

    Price of haircut in India: 300 INR (?) $\sim 3.5$ \$

    But: price of notebooks etc. similar

    \begin{itemize}
        \item tradable vs.~non-tradable goods
        \item commodities vs.~manufactured goods
        \item labour productivity / wages
        \item imbalances / artificial exchange rates
        \item financial account dominates current account?
    \end{itemize}  
    }
\end{frame}

\begin{frame}{Exchange Rates: Interest Rate Parity}
\mytext{
    Investment in both countries yield is equal
    \begin{equation}
        (1+i_\$) = (1+w) \cdot (1+i_\text{\euro} )
    \end{equation}
    with $i$ interest rate,\\ $w$ expected change of exchange rate;

    \medskip

    Fisher equation:
    \begin{equation}
        (1+i) = (1+r) \cdot  (1+\pi)
    \end{equation}
    with $r$ real exchange rate and $\pi$\ inflation
}
    
\end{frame}


\begin{frame}{Labour Supply}
    \showmytikz{../shared/tikz/tikz-laboursupply}
\end{frame}

\begin{frame}{Inflation and NAIRU}
    
\end{frame}

\begin{frame}{IS LM}
    
\end{frame}